\documentclass[conference]{IEEEtran}
\IEEEoverridecommandlockouts
\usepackage{cite}
\usepackage{amsmath,amssymb,amsfonts}
\usepackage{algorithmic}
\usepackage{graphicx}
\usepackage{textcomp}
\usepackage{xcolor}
\usepackage{hyperref}
\usepackage{booktabs}

\def\BibTeX{{\rm B\kern-.05em{\sc i\kern-.025em b}\kern-.08em
    T\kern-.1667em\lower.7ex\hbox{E}\kern-.125emX}}

\begin{document}

\title{A Deterministic, Client-Centric Single Page Architecture for Intercity Bus Ticketing and Session-Resilient Checkout}

\author{\IEEEauthorblockN{Research Engineering Team}
\IEEEauthorblockA{\textit{Department of Computer Science and Engineering} \\
\textit{Antigravity Advanced Computing Lab}\\
Email: research@safar.internal}
}

\maketitle

\begin{abstract}
Intercity bus transit across emerging economies involves millions of daily passenger bookings conducted over variable, high-latency mobile networks. Conventional electronic reservation systems rely on server-rendered round trips and multi-page form reloads, resulting in substantial transaction abandonment when accidental page refreshes flush in-progress booking states. In this paper, we design and implement \textbf{Safar}, a client-centric Single Page Application (SPA) architecture for intercity bus ticketing built with React 19. Safar addresses server contention and network fragility by employing a deterministic procedural synthesis engine that algorithmically derives transit schedules, multi-deck sleeper berth geometry, dynamic pricing, and ladies-priority safety seat allocations without live database dependencies. We introduce a dual-tier client storage model that decouples transient checkout staging in \texttt{sessionStorage} from committed transactional entities in \texttt{localStorage}. Experimental benchmarking demonstrates that the entire application compiles to a minimal footprint of 292.88~kB JavaScript (84.32~kB gzipped) and 44.04~kB CSS, achieving sub-250ms interactive view transitions, zero static code analysis errors across 104 rules, and complete recovery of multi-passenger checkout state across sudden page reloads.
\end{abstract}

\begin{IEEEkeywords}
Bus reservation systems, Single Page Application (SPA), React 19, deterministic procedural generation, client-side state durability, transit ticketing.
\end{IEEEkeywords}

\section{Introduction}
Public road transportation represents the circulatory backbone of regional mobility across India and South Asia. With ubiquitous smartphone adoption, digital ticketing platforms have largely superseded manual bus terminus counters. Despite widespread deployment, commercial bus aggregators frequently suffer from bloated client payloads, slow First Contentful Paint (FCP) latencies, intrusive tracking scripts, and fragile session lifecycles.

When commuters navigate multi-step booking flows on mobile browsers, intermittent connectivity or inadvertent gesture navigations often trigger full page reloads, flushing entered passenger manifests and selected berth coordinates. Concurrently, centralized reservation servers face extreme contention spikes during holiday ticketing surges \cite{pedone2001optimistic}.

To resolve these architectural limitations, this work investigates whether an intercity bus reservation platform can be engineered entirely within a client-centric Single Page Architecture that achieves:
\begin{enumerate}
    \item Complete deterministic inventory generation without live database latency.
    \item Accurate 2D spatial representation of multi-deck sleeper (Lower/Upper) and seater coaches.
    \item Resilient session durability that preserves complex checkout states across reloads.
    \item Fast Time-to-Interactive (TTI) utilizing a zero-overhead CSS token design system.
\end{enumerate}

\section{Related Work}
Object-oriented modeling of transit booking systems has been examined extensively in the software engineering literature. Mohammed and Kassem \cite{mohammed2020uml} presented a UML framework for public bus reservation systems in Egypt, identifying the essential interactions between schedule management, route indexing, and seat assignment. Abu-Dalbouh and Alateyah \cite{abudalbouh2020extension} demonstrated that standard UML diagrams often fail to represent asynchronous web-specific modal transitions and proposed extended stereotypes for web-based transit workflows.

In the domain of web performance, Naeem \cite{naeem2019performance} conducted comparative evaluations between Single Page Applications (SPAs) and Progressive Web Apps (PWAs), concluding that client-rendered SPAs reduce cumulative bandwidth consumption by over 60\% compared to traditional multi-page web applications. In distributed transaction management, Pedone \cite{pedone2001optimistic} formulated optimistic ticket validation protocols to alleviate locking bottlenecks in centralized booking architectures. Safar synthesizes these principles into an actionable, fully verified software implementation.

\section{Problem Statement}
Consider an intercity transit corridor connecting origin city $C_{src}$ to destination $C_{dst}$ on travel date $D$. Traditional booking aggregators require a series of synchronous server queries:
\begin{equation}
    Q_1(C_{src}, C_{dst}, D) \rightarrow \text{BusList}
\end{equation}
\begin{equation}
    Q_2(\text{BusID}) \rightarrow \text{SeatMatrix}
\end{equation}
\begin{equation}
    Q_3(\text{UserID}, \text{CouponCode}) \rightarrow \text{Discount}
\end{equation}
\begin{equation}
    Q_4(\text{BookingPayload}) \rightarrow \text{Ticket}
\end{equation}

Under congested network conditions, failure or latency during any intermediate query $Q_k$ invalidates previous passenger selections. Furthermore, aggregators frequently fail to convey physical coach geometry, particularly the distinction between Lower and Upper berths in sleeper coaches and designated female-passenger priority allocations. The problem addressed by this paper is formulating an architecture that executes the discovery, selection, validation, and staging pipeline entirely client-side while guaranteeing mathematical reproducibility and session durability.

\section{Proposed Methodology}

\subsection{Deterministic Procedural Data Synthesis}
To achieve zero-latency inventory queries without hardcoded databases, Safar derives all transit entities from a deterministic integer seed computed from the search criteria tuple:
\begin{equation}
\text{Seed} = \sum_{k=0}^{|S|-1} \text{charCodeAt}(S[k])
\end{equation}
Where $S = C_{src} \parallel \text{"-"} \parallel C_{dst} \parallel \text{"-"} \parallel D$. Using modular arithmetic seeded by $\text{Seed}$, the engine computes:
\begin{itemize}
    \item Service frequency: $N = 5 + (\text{Seed} \pmod 4) \in [5, 8]$.
    \item Departure timeline: hours selected from $\mathcal{H} = \{06, 08, 14, 20, 21, 22, 23\}$.
    \item Geographic duration: base transit duration scaled by route distance and seeded jitter ($\pm 20 \text{ minutes}$).
    \item Dynamic pricing: base fare scaled by coach type modifier ($1.40\times$ for A/C Sleeper, $1.15\times$ for A/C Seater, $0.85\times$ for Non-A/C).
\end{itemize}

\subsection{Geometric Multi-Deck Seat Allocation}
Sleeper coaches represent a 3-dimensional physical arrangement mapped onto a 2D interactive grid:
\begin{itemize}
    \item Lower Deck ($\mathcal{D}_L$) and Upper Deck ($\mathcal{D}_U$), each containing 15 berths ($|\mathcal{D}| = 30$).
    \item Berths follow a transverse 1+2 spatial layout: Left Window ($L_w$), Right Aisle ($R_a$), Right Window ($R_w$).
    \item Deterministic occupancy and safety predicates:
    \begin{equation}
        \text{Booked}(i) \iff (\text{Seed} + 11i) \pmod 6 = 0
    \end{equation}
    \begin{equation}
        \text{LadiesPriority}(i) \iff (\text{Seed} + 17i) \pmod{13} = 0 \quad (\text{if } \neg\text{Booked}(i))
    \end{equation}
\end{itemize}
The choice of coprime primes ($11, 17$) guarantees uniform scattering across rows and columns without spatial clustering.

\subsection{Dual-Tier State Coordination Model}
To prevent form-state destruction during mid-booking authentication gates, state is bifurcated:
\begin{itemize}
    \item \textbf{Ephemeral Buffer (\texttt{sessionStorage})}: Stores \texttt{safar\_checkout\_bus}, \texttt{safar\_selected\_seats}, and \texttt{safar\_checkout\_form} (passenger names, ages, contact information).
    \item \textbf{Persistent Relational Store (\texttt{localStorage})}: Stores \texttt{safar\_users} and \texttt{safar\_bookings}. Mutations commit atomically upon payment approval.
\end{itemize}

\section{System Architecture}
The system architecture decomposes into four decoupled tiers:
\begin{enumerate}
    \item \textbf{Presentation Tier}: React 19 virtual DOM, custom hash-based router, and 13 reusable UI component primitives.
    \item \textbf{State Coordination Tier}: React hooks coupled with \texttt{sessionStorage} for transient checkout staging.
    \item \textbf{Business Logic Engine}: Algorithmic bus and seat matrix synthesis, fare calculation, and promo voucher verification.
    \item \textbf{Simulated Relational Persistence}: \texttt{mockDb} engine executing atomic CRUD transactions against \texttt{localStorage}.
\end{enumerate}

\section{Implementation Details}
The platform is developed in ECMAScript 2022 and bundled using Vite v8.2.2.
\begin{itemize}
    \item \textbf{Design System}: Encapsulated in \texttt{src/index.css} utilizing native CSS custom properties (\texttt{--brand-primary: \#1e3a8a}, \texttt{--brand-accent: \#0284c7}).
    \item \textbf{Promotional Verification}: Functional predicates validate voucher rules:
    \begin{itemize}
        \item \texttt{FIRSTTRIP}: 15\% discount (max ₹150) restricted to users with zero prior bookings.
        \item \texttt{ROUTE10}: Flat ₹100 deduction for fares $\ge ₹600$.
        \item \texttt{WEEKEND}: 10\% discount (max ₹250) valid only on Saturdays and Sundays.
    \end{itemize}
    \item \textbf{Print Optimization}: Dedicated \texttt{@media print} stylesheets isolate the ticket component while stripping browser chrome.
\end{itemize}

\section{Experimental Setup \& Results}
The implementation was validated on an x86\_64 multi-core workstation running Node.js v23.10.0:
\begin{itemize}
    \item \textbf{Static Code Analysis}: Oxlint v1.79.0 evaluated 104 rules across all 30 source files in 137ms, yielding \textbf{0 errors and 0 warnings}.
    \item \textbf{Production Compilation}: Vite v8.2.2 completed production bundling in \textbf{530ms}.
    \begin{table}[htbp]
    \caption{Production Bundle Distribution Metrics}
    \begin{center}
    \begin{tabular}{lrr}
    \toprule
    \textbf{Asset Type} & \textbf{Uncompressed} & \textbf{Gzip Compressed} \\
    \midrule
    \texttt{index.html} & 0.66 kB & 0.41 kB \\
    \texttt{index.css} & 44.04 kB & 8.60 kB \\
    \texttt{index.js} & 292.88 kB & 84.32 kB \\
    \textbf{Total Assets} & \textbf{337.58 kB} & \textbf{93.33 kB} \\
    \bottomrule
    \end{tabular}
    \end{center}
    \end{table}
    \item \textbf{User Journey Execution}: An autonomous browser subagent traversed the full booking lifecycle from search to ticket issuance (\texttt{BK-662345}), verifying sub-250ms view transitions and complete state preservation across reloads.
\end{itemize}

\section{Limitations}
Data stored in \texttt{localStorage} is scoped to the individual client browser profile and does not synchronize across separate physical devices. Furthermore, payment transactions are validated syntactically but simulated without real fund capture.

\section{Conclusion and Future Work}
This paper presented the design, implementation, and experimental evaluation of Safar, a client-centric Single Page Architecture for intercity bus ticketing. By combining deterministic procedural generation with a dual-tier client storage model, Safar eliminates server latency, preserves form state across reloads, and visualizes multi-deck seat matrices cleanly. Future work will investigate WebSocket telemetry for live vehicle tracking and service worker caching for offline ticket validation.

\bibliographystyle{IEEEtran}
\bibliography{references}

\end{document}
